\medskip

La solution générale d'une équation différentielle d'ordre 1 est la somme d'une solution particulière et de la solution générale de l'équation différentielle sans second membre associée.

\begin{list}{\textbullet}{}
\item L'équation sans second membre associée est: $y'=-3y$.

D'après le cours, on sait que les solutions de cette équation sont les fonctions $y_0$ définies par $y_0(x)= k\e^{-3x}$ où $k$ est un réel quelconque.

\item On cherche une solution particulière de $(E)$ sous forme de constante $a$. \\
Donc $a'=-3a+5$. Or $a'=0$, donc $-3a+5=0$ donc $a=\dfrac{5}{3}$.
\end{list}

Les fonctions solutions de l'équation différentielle $(E)$ sont les fonctions $y$ définies sur $\R$ par $y(x)= k \e^{-3x} +\dfrac{5}{3}$ où $k$ est un réel quelconque.

\bigskip
